
\section{Experimental Evaluation}\label{sec:experiment}

The focus of these experiments is to evaluate \model in terms of its ability to enable effective, generalizable decision making. In particular, we evaluate
\begin{itemize}[leftmargin=2em, itemsep=0pt]
    \item[{\bf (1)}] the ability to combinatorially generalize across different subgoals in \sect{sect:combinatorical}, 
    \item[{\bf (2)}] the ability to effectively learn and generalize across many tasks in \sect{sect:multi},
    \item[{\bf (3)}] the ability to leverage existing videos on the internet to generalize to complex tasks in \sect{sect:real}.
\end{itemize}
See experimental details in Appendix~\ref{app:details}. 
Additional results are given in Appendix~\ref{app:results} and videos 
in the supplement.


\subsection{Combinatorial Policy Synthesis}
\label{sect:combinatorical}

First, we measure the ability of \model to combinatorially generalize to different language tasks.

\myparagraph{Setup.} To measure combinatorial generalization, we use the combinatorial robot planning tasks in \citep{Mao2022PDSketch}. In this task, a robot must manipulate blocks in an environment to satisfy language instructions, i.e., {\it put a red block right of a cyan block}. To accomplish this task, the robot must first pick up a white block, place it in the appropriate bowl to paint it a particular color, and then pick up and place the block in a plate so that it satisfies the specified relation. In contrast to \citep{Mao2022PDSketch} which uses pre-programmed pick and place primitives for action prediction, we predict actions in the continuous robotic joint space for both the baselines and our approach.

We split the language instructions in this environment into two sets: one set of instructions (70\%) that is seen during training, and another set (30\%) that is only seen during testing. The precise locations of individual blocks, bowls, and plates in the environment are fully randomized in each environment iteration. We train the video model on 200k example videos of generated language instructions in the train set. Details of this environment can be found in Appendix~\ref{app:details}. We constructed demonstrations of videos in this task by using a scripted agent.

\input{tbls/combinatorical.tex}
\input{figText/combinatorical_plan.tex}

\input{figText/action_execution} 
\vspace{-7pt}
\paragraph{Baselines.} We compare the proposed approach with three separate representative approaches. First, we compare to existing work that uses goal-conditioned transformers to learn across multiple environments, where goals can be specified as episode returns \citep{lee2022multi}, expert demonstrations \citep{reed2022generalist}, or text and images \citep{brohan2022rt}. To represent these baselines, we construct a transformer behavior cloning (BC) agent to predict the subsequent action to execute given the task description and either the visual observation (Image + Transformer BC) or the underlying robot joint state (State + Transformer BC). Second, given that our approach regresses a sequence of actions to execute, we further compare with transformer models that regress a sequence of future actions to execute, similar to the goal-conditioned behavioral cloning of the Trajectory Transformer~\citep{janner2021offline} (Image + TT). Finally, to highlight the importance of the video-as-policy approach, we compare \model with learning a diffusion process that, conditioned on an image observation, directly infers future robot actions in the joint space (as opposed to diffusing future image frames), corresponding to \citep{janner2022planning,ajay2022conditional}. For both our method and each baseline, we condition the policy on encoded language instructions using pretrained T5 embeddings. 
 Note that in this setting, existing offline reinforcement learning baselines are not directly applicable as we do not have access to the reward functions in the environment.
 
\input{figText/multitask_plan.tex}
\input{tbls/ablations}
\myparagraph{Metrics.} To compare \model with baselines, we measure final task completion accuracy across new instances of the environment and associated language prompts. We subdivide the evaluation along two axes: \textbf{(1)} whether the language instruction has been seen during training and \textbf{(2)} whether the language instruction specifies placing a block in relation to some other block as opposed to direct pick-and-place.

\input{figText/real_plan.tex}

\myparagraph{Combinatorial Generalization.} In \tbl{tbl:combinatorical}, we find that \model generalizes well to both seen and novel combinations of language prompts . We illustrate our action generation pipeline in \fig{fig:comb_pipeline} and different generated video plans using our approach in \fig{fig:comb_qual}.

\input{figText/adaptability}
\vspace{-7pt}
\looseness=-1
\paragraph{Ablations.} In \tbl{tbl:ablations}, we ablate \model on seen language instructions and in-relation-to tasks. Specifically, we study the effect of conditioning the video generative model on the first observation frame (frame condition), tiling the observed frame across timesteps (frame consistency) and super-resolving video generation across time (temporal hierarchy). In settings where frame consistency is not enforced, we provide a zeroed out image as context to the non-start frames in a video. We found that all components of \model are crucial for good performance. We found that frame conditioning and consistency enabled videos to be consistent with the observed image and temporal hierarchy enabled more detailed video plans. 


\myparagraph{Adaptability.} We next assess the ability of \model to adapt at test time to new constraints.  In \fig{fig:comb_adapt}, we illustrate the ability to construct plans which color and move one particular block to a specified geometric relation. 

\subsection{Multi-Environment Transfer}
\label{sect:multi}
\input{tbls/multitask.tex}



We next evaluate the ability of \model to effectively learn across a set of different tasks and generalize, at test time, to a new set of unseen environments.

\myparagraph{Setup.} To measure multi-task learning and transfer, we use the suite of language guided manipulation tasks from \citep{shridhar2022cliport}. We train our method using demonstrations across a set of 10 separate tasks from \citep{shridhar2022cliport}, and evaluate the ability of our approach to transfer to 3 different test tasks. 
Using a scripted oracle agent, we generate a set of 200k videos of language execution in the environment. We report the underlying accuracy in which each language instruction is completed. Details of this environment can be found in Appendix~\ref{app:details}.

\myparagraph{Baselines.} We use the same baseline methods as in \sect{sect:combinatorical}. 
While our environment setting is similar to that of \citep{shridhar2022cliport}, this method is not directly comparable to our approach, as CLIPort abstracts actions to the existing primitives of pick and place as opposed to using the joint space of a robot. CLIPort is also designed to solve the significantly simpler problem of inferring only the poses upon which to pick and place objects (with no easy manner to adapt to our setting). 

\myparagraph{Multitask Generalization.} In \tbl{tbl:multitask} we present results of \model %
and baselines across new tasks. %
The \model approach is able to generalize and synthesize new videos and decisions of different language tasks, and can generate videos consisting of picking different kinds of objects and different colored objects. We further present video visualizations of our approach in \fig{fig:multitask_qual}.

\subsection{Real World Transfer}
\label{sect:real}


Finally we evaluate the extent to which \model can generalize to real world scenarios and construct complex behaviors by leveraging widely available videos on the internet.

\myparagraph{Setup.} Our training data consists of an internet-scale pretraining dataset and a smaller real-world robotic dataset. The pretraining dataset uses the same data as \citep{ho2022imagen}, which consists of 14 million video-text pairs, 60 million image-text pairs, and the publicly available LAION-400M image-text dataset. The robotic dataset is adopted from the Bridge dataset~\citep{ebert2021bridge} with 7.2k video-text pairs, where we use the task IDs as texts. We partition the 7.2k video-text pairs into train (80\%) and test (20\%) splits. We pretrain \model on the pretraining dataset followed by finetuning on the train split of the Bridge data. Architectural details can be found in Appendix~\ref{app:details}. 

\input{figText/real_generalization_ood.tex}


\myparagraph{Video Synthesis.} We are particularly interested in the effect of pretraining on internet-scale video data that is not specific to robotics. We report the CLIP scores, FIDs, and FVDs (averaged across frames and computed on 32 samples) of \model trained on Bridge data, with and without pretraining. As shown in \tbl{tbl:real}, \model with pretraining achieves significantly higher FID and FVD and a marginally better CLIP score than \model without pretraining, suggesting that pretraining on non-robot data helps with generating plans for robots. Interestingly, \model without pretraining often synthesizes plans that fail to complete the task (\fig{fig:real_qual}), which is not well reflected in the CLIP score, suggesting the need for better generation metrics for control-specific tasks. To tackle the lack of such a metric, we develop a surrogate metric for evaluating task success from the generated videos. Specifically, we train a success classifier that takes in the last frame of a generated video and predicts whether the task is successful or not. We find that training a model from scratch achieves 72.6\% success while finetuning from a pretrained model improves performance to 77.1\%. In both settings, generated videos are able to successfully complete most tasks.

\myparagraph{Generalization.} We find that internet-scale pretraining enables \model to generalize to novel task commands and scenes in the test split not seen during training, whereas \model trained only on task-specific robot data fails to generalize. Specifically, \fig{fig:real_plan_generalize} shows the %
results of novel task commands that do not exist in the Bridge dataset. Additionally, \model is relatively robust to background changes such as black cropping or the addition of photo-shopped objects as shown in \fig{fig:real_ood}.


\input{tbls/real_world.tex}